\begin{table}[!htbp]
\centering
\caption{Victimization categories and cumulative treatment coverage}
\label{tab:desc_victimization_categories}
\begin{tabular}{lrrrrrrrr}
\toprule
Category & N & Score min. & Score median & Score max. & Treated 2013 (\%) & Treated 2017 (\%) & Treated 2023 (\%) & Wellbeing 2007 \\
\midrule
A & 1,284 & 0.1538 & 0.2898 & 2.3622 & 59.7 & 85.2 & 91.3 & 0.304 \\
B & 1,269 & 0.0623 & 0.0989 & 0.1537 & 42.0 & 61.1 & 86.0 & 0.309 \\
C & 1,310 & 0.0269 & 0.0406 & 0.0623 & 27.4 & 27.9 & 84.0 & 0.330 \\
D & 1,127 & 0.0152 & 0.0197 & 0.0300 & 13.3 & 13.7 & 63.5 & 0.340 \\
E & 722 & 0.0077 & 0.0102 & 0.0150 & 2.9 & 2.9 & 19.5 & 0.338 \\
\bottomrule
\end{tabular}
\parbox{0.98\linewidth}{\footnotesize \textit{Notes:} Categories A-E and the observed score are supplied by RUV. Treatment percentages are unconditional means of the cumulative CMAN indicators. Baseline wellbeing is the mean equal-domain 2007 score among linked communities with complete components. The rounded RUV score can tie across official six-decimal thresholds; this table does not infer assignment side from the score or conduct adjacent-category tests. Sources: RUV, CMAN 2023, and 2007 Census CCPP tabulation.}
\end{table}
